\section{The difficulty atlas}
\label{sec:atlas}

The estimator returns two kinds of object. One is per-rider: the skill posteriors
$\hat\skill_\rider$. The other is a property of the environment: the difficulty
function $\hat\diff(\metric,\site)$ itself. Collected across sites and dimensions,
the latter forms what we call the \emph{intrinsic difficulty atlas}---a
site-level, anonymous map of how hard each condition is, net of who rides it.

\begin{figure}[t]
  \centering
  \includegraphics[width=0.86\linewidth]{difficulty-atlas.pdf}
  \caption{The construct the atlas measures. \emph{Bottom:} the intrinsic
  difficulty function $\diff(\metric)$ (dashed), minimised at the ground-truth
  optimum $\metric^\star = 18$~kn; the recovered $\hat\diff(\metric)$ is overlaid
  where available. \emph{Top:} three skill-conditioned suitability curves
  $\sigmoid(a(\skill - \diff(\metric)))$ for beginner ($\skill{=}{-}1$),
  intermediate ($\skill{=}0$), and expert ($\skill{=}{+}1$). Skill shifts the
  suitability curve while the difficulty beneath it is invariant: two operators
  whose spots share the same $\diff$ but differ in clientele read different
  suitability surfaces from the \emph{same} underlying difficulty. That
  separation is what no meteorological observation network records. Figure
  generated from code (Appendix~\ref{app:repro}).}
  \label{fig:atlas}
\end{figure}

\subsection{Why the atlas is new}

Existing observation networks---marine reanalyses, atmospheric mesonets,
lightning-detection and air-quality networks, national avalanche services---
densely cover the \emph{meteorological} state. What none of them holds is the
quantity $\diff$ isolates: the intrinsic difficulty of a condition for a skilled
human activity, separated from the population skill mix that otherwise confounds
it. Suitability, as conventionally fit, is a blend of difficulty and clientele;
$\diff$ is the difficulty term alone. Because it is anchored to a physical curve
(Section~\ref{sec:model}), it is expressed in interpretable units and is
\emph{joinable} to those same meteorological archives: for any site and condition,
the atlas adds a difficulty coordinate the archives lack.

\subsection{A measurable construct}

The consequence is conceptual as much as practical. Two spots with identical
fitted suitability curves can have very different intrinsic difficulty, masked by
different clienteles; until the skill term is identified, they are
indistinguishable. Inverse Suitability makes them distinguishable, and in doing so
treats activity suitability---for the first time, to our knowledge---as a formal,
measurable latent construct rather than an expert convention. The atlas is the
artifact of that construct: a difficulty rating decoupled from who happens to be
present, defined precisely enough to be estimated, reproduced, and audited.
